\subsection{Exploratory actual-action benefit}
After inspecting the fixed study, we reanalyzed all saved predictions to measure
the signed benefit of each actual sparse graph, rather than only full expansion.
Previous-risk correlation with its own 25\% action's normalized vector benefit is
$-0.021\pm0.049$ for faithful and $-0.052\pm0.013$ for NLL.
The corresponding dense and actual-risk benefit signs disagree for
$30.2\pm0.6$\% and
$31.3\pm0.5$\% of particles, respectively.
Thus dense benefit is an imperfect label for the sparse action, but replacing
it with actual action benefit does not recover strong risk ranking here.
Weak average rank association does not exclude nonmonotone structure: the highest
risk quartile has positive mean risk-action benefit for each objective, although
one NLL seed is negative and random allocation yields larger benefit in that
quartile in all six models.

A hindsight whole-frame choice among base, random25, speed25 and risk25 selects
base as the strictly best action on $35.9\pm12.0$\% / $35.8\pm6.3$\% of frames (faithful/NLL).
This target-using diagnostic motivates studying when to add edges; it is not a
deployable selector or a free-rollout result. All values retain equal-trajectory
weighting and sample SD across three seeds. These post-inspection findings remain
conditional on the original evaluation graph convention.

\begin{figure}[t]
\centering\setlength{\unitlength}{1pt}
\begin{picture}(225,167)
\put(12,155){\makebox(0,0)[l]{\scriptsize Previous-risk Spearman correlation}}
\put(30,31){\line(1,0){181}}
\put(30,31){\line(0,1){114}}
\put(27,31.00){\makebox(0,0)[r]{\tiny -0.1}}
\put(29,31.00){\line(1,0){3}}
\put(27,49.00){\makebox(0,0)[r]{\tiny 0.0}}
\put(29,49.00){\line(1,0){3}}
\put(27,67.00){\makebox(0,0)[r]{\tiny 0.1}}
\put(29,67.00){\line(1,0){3}}
\put(27,85.00){\makebox(0,0)[r]{\tiny 0.2}}
\put(29,85.00){\line(1,0){3}}
\put(27,103.00){\makebox(0,0)[r]{\tiny 0.3}}
\put(29,103.00){\line(1,0){3}}
\put(27,121.00){\makebox(0,0)[r]{\tiny 0.4}}
\put(29,121.00){\line(1,0){3}}
\put(27,139.00){\makebox(0,0)[r]{\tiny 0.5}}
\put(29,139.00){\line(1,0){3}}
\multiput(32,49)(4,0){45}{\line(1,0){1}}
\put(50.00,101.43){\circle*{3.5}}
\put(52.00,122.47){\circle*{3.5}}
\put(54.00,115.92){\circle*{3.5}}
\put(48.00,113.27){\line(1,0){8}}
\put(111.00,45.20){\circle*{3.5}}
\put(113.00,53.06){\circle*{3.5}}
\put(115.00,36.93){\circle*{3.5}}
\put(109.00,45.06){\line(1,0){8}}
\put(172.00,47.70){\circle*{3.5}}
\put(174.00,52.41){\circle*{3.5}}
\put(176.00,35.43){\circle*{3.5}}
\put(170.00,45.18){\line(1,0){8}}
\put(64.00,120.05){\makebox(0,0){\tiny$\triangle$}}
\put(66.00,127.68){\makebox(0,0){\tiny$\triangle$}}
\put(68.00,121.38){\makebox(0,0){\tiny$\triangle$}}
\put(62.00,123.04){\line(1,0){8}}
\put(125.00,42.54){\makebox(0,0){\tiny$\triangle$}}
\put(127.00,38.39){\makebox(0,0){\tiny$\triangle$}}
\put(129.00,40.05){\makebox(0,0){\tiny$\triangle$}}
\put(123.00,40.33){\line(1,0){8}}
\put(186.00,42.11){\makebox(0,0){\tiny$\triangle$}}
\put(188.00,39.52){\makebox(0,0){\tiny$\triangle$}}
\put(190.00,37.41){\makebox(0,0){\tiny$\triangle$}}
\put(184.00,39.68){\line(1,0){8}}
\put(59,21){\makebox(0,0){\tiny Base residual}}
\put(120,21){\makebox(0,0){\tiny Dense benefit}}
\put(181,21){\makebox(0,0){\tiny Actual risk25}}
\put(68,7){\circle*{3.5}}
\put(74,7){\makebox(0,0)[l]{\tiny Faithful}}
\put(136,7){\makebox(0,0){\tiny$\triangle$}}
\put(143,7){\makebox(0,0)[l]{\tiny NLL}}
\end{picture}
\caption{Residual ranking and actual action value on the saved observed histories. Each symbol is one seed mean; horizontal marks are three-seed means. Actual risk25 benefit uses the whole selected graph. This exploratory comparison retains the original no-self-loop evaluation convention.}
\label{fig:actual-action-risk}
\end{figure}
